\section{Scientific Artifact Architecture}
\subsection{Experiment manifests}
Every planned run is generated from a machine-readable condition containing dataset and
version, architecture, training seed, deletion strategy and fraction, method and
hyperparameters, exact seeds, metrics, and expected hardware class.  Reduced, main, and
extended profiles make the size of the experiment explicit before execution.  Expansion
calculates the number of trainings, approximate device-hours, checkpoint volume, and result
count, preventing an accidental factorial explosion.

\subsection{Canonical data lineage}
The lineage chain is
\begin{center}
source dataset $\rightarrow$ preprocessing $\rightarrow$ canonical split
$\rightarrow$ deletion manifest $\rightarrow$ run configuration
$\rightarrow$ checkpoint and predictions $\rightarrow$ metrics
$\rightarrow$ statistics $\rightarrow$ figure or table $\rightarrow$ claim.
\end{center}
Each edge carries a stable identifier or cryptographic hash.  Split indices and deletion
indices are archived rather than reconstructed from prose.  Dataset validation records
schema, counts, class balance, duplicate rate, missing values, label validity, overlap, and
source integrity in both human- and machine-readable forms.

Lineage is deliberately more granular than an experiment name. A result that only records
``Fashion-MNIST, seed 11'' cannot establish which preprocessing revision, split, or
deletion request created it. The stored chain makes counterfactual exclusion auditable:
the exact-training manifest can be compared directly with the forgotten-ID manifest rather
than inferred from control flow. This also permits a later reviewer to distinguish a model
trained under a new preprocessing version from a genuine rerun of an earlier condition.

\subsection{Append-only raw records}
A completed run receives a unique identifier, timestamp, configuration hash, Git commit,
environment snapshot, device, runtime, seed, metrics, and checkpoint references.  A rerun
creates a new record and may point to its parent; it never replaces the earlier payload.
Raw predictions remain available for metric recomputation.  Processed summaries are
rebuildable and are kept physically distinct from raw evidence.

Immutability is not a claim that files are technologically impossible to alter; it is an
operational contract. New execution creates new evidence, and validation records the
configuration and content identifiers required to detect substitution or accidental
overwrite. Derived artifacts may be regenerated when analysis code changes, but the
derivation must identify the raw inputs and software revision used. This separation allows
an analysis bug to be corrected without rewriting the history of the run that exposed it.

\subsection{Acceptance gates}
A run enters the paper dataset only when dataset and split hashes match, retained and
forgotten partitions are disjoint, every forgotten identifier is absent from exact
training, training completed, predictions exist, metrics validate, metadata is complete,
and unexplained non-finite values are absent.  Gate success means that these properties
were checked; it does not mean that the unlearning method succeeded.  Failed runs remain in
a separate registry with category, exception, configuration, batch impact, and remediation.

The acceptance gate is intentionally conservative. It cannot prove that a method achieved
unlearning, but it can reject records whose numerical values would otherwise invite an
unsupported interpretation. For example, a completed training job is insufficient if its
forgotten examples leaked into the exact-retraining dataset, if predictions for a required
partition are absent, or if a metric silently contains non-finite values. Failed runs are
not pooled as observations, yet their presence matters for operational reproducibility and
for assessing whether a stated experimental matrix was actually completed.

\subsection{Claim--evidence graph}
Substantive quantitative claims are represented as nodes that point to experiment IDs,
processed results, statistical analyses, and figure or table outputs.  Release validation
rejects a verified claim if any referenced artifact is absent or hash-inconsistent.  This
mechanism reduces manual transcription and distinguishes measured statements from
interpretation, hypotheses, and limitations.  Headline LaTeX numbers and tables are
generated from validated artifacts rather than copied into prose.

\subsection{Release and independent reproduction}
The release bundle contains source, locked configuration, environment records, immutable
evidence, analysis programs, generated values, figures, tables, manuscript source, and
checksums.  Level-A reproduction starts from archived predictions; Level-B reproduction
repeats acquisition and training; Level-C replication independently implements the written
protocol.  These levels are reported separately because a reproducible script does not by
itself demonstrate independent scientific replication.

\subsection{Private development disclosure record}
Substantial automated assistance is recorded privately by activity class---implementation,
orchestration, analysis, visualization, and drafting---so that the human author can comply
with a venue's disclosure policy.  Prompt transcripts and irrelevant scratch traces are not
scientific provenance and need not be public.  Dataset lineage, executed configurations,
failures, numerical evidence, and manuscript claim links are scientific provenance and are
retained regardless of the development tools used.
